\documentclass[10pt,a4paper]{amsart}
\usepackage{amsmath,amssymb,amsthm,graphicx}
% Retained minimal preamble for the Stacks Project chapter more-algebra.tex (More on Algebra).
% The chapter file pinned-original-source.tex carries only \input{preamble} on its first line, so no
% preamble text lives inside the chapter itself.  The Stacks Project uses one single shared preamble for
% every chapter; the snapshot stages it identically as sources/stacks-chow/preamble.tex and
% sources/stacks-etale-cohomology/preamble.tex.  The block below is copied mechanically, line for line,
% from sources/stacks-chow/preamble.tex lines 14-18 (the Xy-pic setup) and 154-217 (the native theorem
% declarations together with the mathematical macros).  Nothing was added, removed or redefined; the
% wrapper class and the amsmath/amssymb/amsthm/graphicx base packages are supplied by the assembler.
% Chapter: more-algebra.tex -- More on Algebra.
% For commutative diagrams we use Xy-pic
\usepackage[all]{xy}

% We use 2cell for 2-commutative diagrams.
\xyoption{2cell}

% Theorem environments.
%
\theoremstyle{plain}
\newtheorem{theorem}[subsection]{Theorem}
\newtheorem{proposition}[subsection]{Proposition}
\newtheorem{lemma}[subsection]{Lemma}

\theoremstyle{definition}
\newtheorem{definition}[subsection]{Definition}
\newtheorem{example}[subsection]{Example}
\newtheorem{exercise}[subsection]{Exercise}
\newtheorem{situation}[subsection]{Situation}

\theoremstyle{remark}
\newtheorem{remark}[subsection]{Remark}
\newtheorem{remarks}[subsection]{Remarks}

\numberwithin{equation}{subsection}

% Macros
%
\def\lim{\mathop{\mathrm{lim}}\nolimits}
\def\colim{\mathop{\mathrm{colim}}\nolimits}
\def\Spec{\mathop{\mathrm{Spec}}}
\def\Hom{\mathop{\mathrm{Hom}}\nolimits}
\def\Ext{\mathop{\mathrm{Ext}}\nolimits}
\def\SheafHom{\mathop{\mathcal{H}\!\mathit{om}}\nolimits}
\def\SheafExt{\mathop{\mathcal{E}\!\mathit{xt}}\nolimits}
\def\Sch{\mathit{Sch}}
\def\Mor{\mathop{\mathrm{Mor}}\nolimits}
\def\Ob{\mathop{\mathrm{Ob}}\nolimits}
\def\Sh{\mathop{\mathit{Sh}}\nolimits}
\def\NL{\mathop{N\!L}\nolimits}
\def\CH{\mathop{\mathrm{CH}}\nolimits}
\def\proetale{{pro\text{-}\acute{e}tale}}
\def\etale{{\acute{e}tale}}
\def\QCoh{\mathit{QCoh}}
\def\Ker{\mathop{\mathrm{Ker}}}
\def\Im{\mathop{\mathrm{Im}}}
\def\Coker{\mathop{\mathrm{Coker}}}
\def\Coim{\mathop{\mathrm{Coim}}}

% Boxtimes
%
\DeclareMathSymbol{\boxtimes}{\mathbin}{AMSa}{"02}

%
% Macros for moduli stacks/spaces
%
\def\QCohstack{\mathcal{QC}\!\mathit{oh}}
\def\Cohstack{\mathcal{C}\!\mathit{oh}}
\def\Spacesstack{\mathcal{S}\!\mathit{paces}}
\def\Quotfunctor{\mathrm{Quot}}
\def\Hilbfunctor{\mathrm{Hilb}}
\def\Curvesstack{\mathcal{C}\!\mathit{urves}}
\def\Polarizedstack{\mathcal{P}\!\mathit{olarized}}
\def\Complexesstack{\mathcal{C}\!\mathit{omplexes}}
% \Pic is the operator that assigns to X its picard group, usage \Pic(X)
% \Picardstack_{X/B} denotes the Picard stack of X over B
% \Picardfunctor_{X/B} denotes the Picard functor of X over B
\def\Pic{\mathop{\mathrm{Pic}}\nolimits}
\def\Picardstack{\mathcal{P}\!\mathit{ic}}
\def\Picardfunctor{\mathrm{Pic}}
\def\Deformationcategory{\mathcal{D}\!\mathit{ef}}
\usepackage{cleveref}
\crefname{originallemma}{original lemma}{original lemmas}
\Crefname{originallemma}{Original lemma}{Original lemmas}
\makeatletter
\newlabel{more-algebra-lemma-exact-category-stably-free}{{0BC4}{source}{}{sourcepointer-1}{}}
\newlabel{more-algebra-lemma-exact-category-stably-free@cref}{{[originallemma][0][]0BC4}{[1][0][]source}{}{}{}}
\makeatother
\title{In a short exact sequence of finite projective modules, stable freeness of two terms forces it for the third}
\author{The Stacks Project Authors}
\date{stacks-project:more-algebra.tex:local-snapshot-20261012; GNU FDL 1.2 or later}
\begin{document}
\maketitle
\setcounter{subsection}{1}
\noindent\textit{Editorial note.} All delivered mathematical text is the author's own, copied from the pinned Stacks Project chapter file more-algebra.tex; no statement or proof step was rewritten, reordered or removed.\par\medskip
\noindent\textit{Editorial note.} The chapter file pinned-original-source.tex carries only \textbackslash{}input\{preamble\} on its first line, so no preamble text lives inside the chapter. The retained file sources/stacks-more-algebra/retained-preamble.tex is a mechanical, line-for-line copy of the project's single shared preamble as staged in the snapshot, namely sources/stacks-chow/preamble.tex lines 14-18 (Xy-pic setup) and 154-217 (theorem declarations and mathematical macros). Nothing was added, removed or redefined, no package was added, and the excerpt serves only as preamble, never as a span coordinate.\par\medskip
\noindent\textit{Editorial note.} The statement span starts at the \textbackslash{}begin\{lemma\} line and the proof span at the \textbackslash{}begin\{proof\} line, so the two environments are delivered whole; the intermediate definition of stably free modules (source lines 55-64) and the one-line remark at source line 66-67 are deliberately outside the counted spans and are carried as standard prerequisites instead.\par\medskip
\noindent\textit{Editorial note.} The source URL is the source-level Stacks Project URL, as required when more than one card is produced from one source; the per-claim canonical tag 0BC4 is carried by claim\_key and by the external reference. The citation fields of this card (source\_id, source\_version, authors, source\_url, license.name) are identical to those of the two cards already published from this same chapter, and source\_title repeats their section-level wording verbatim, so that the batch selector treats the three as one source instead of deferring them as a same-source clash.\par\medskip
\noindent\textit{Allowed standard prerequisites (editorial).}\begin{itemize}
\item A short exact sequence of projective modules splits. used at proof line 77, from the projectivity of the three modules recorded in the statement at source line 71-72. Turns the sequence 0 -> P' -> P -> P'' -> 0 into an isomorphism P = P' + P'', which is the decomposition all three cases are computed in.
\item Stably free and stably isomorphic modules, as defined in the preceding definition of the same section. definition environment at source lines 55-64 of the same chapter, whose items (1) and (2) define stable isomorphism and stable freeness. Supplies the meaning of the hypothesis 'two out of three are stably free' and of the conclusion 'so is the third', namely freeness after adding a suitable free direct summand.
\item Adding a free direct summand preserves freeness, and direct sums of free modules are free. used at proof lines 79, 84-87 and 88-89 when the free modules R\textasciicircum{}\{n+m\}, R\textasciicircum{}m and R\textasciicircum{}n are assembled. Realizes the stable isomorphism explicitly: it is what lets a module M with M + R\textasciicircum{}n free be recognised as stably free, and it is applied under the regrouping of the direct sum in the second case.
\end{itemize}\medskip
\ifdefined\hypertarget\hypertarget{sourcepointer-1}{}\fi
\noindent\textit{Source reference.} Original lemma 0BC4: focused-seeds/AGP/sources/stacks-snapshot-20261011/tags.txt, source lines 14267--14267. The lemma's own canonical tag, quoted so that readers can resolve the label inside the delivered text; it points at the counted statement itself and carries no proof content.\par\medskip
\begin{lemma}
\label{lemma-exact-category-stably-free}
Let $R$ be a ring. Let $0 \to P' \to P \to P'' \to 0$ be a short
exact sequence of finite projective $R$-modules. If $2$ out of $3$
of these modules are stably free, then so is the third.
\end{lemma}

\begin{proof}
Since the modules are projective, the sequence is split. Thus we can
choose an isomorphism $P = P' \oplus P''$. If $P' \oplus R^{\oplus n}$
and $P'' \oplus R^{\oplus m}$ are free, then we see that
$P \oplus R^{\oplus n + m}$ is free. Suppose that $P'$ and $P$ are
stably free, say $P \oplus R^{\oplus n}$ is free and $P' \oplus R^{\oplus m}$
is free. Then
$$
P'' \oplus (P' \oplus R^{\oplus m}) \oplus R^{\oplus n} =
(P'' \oplus P') \oplus R^{\oplus m} \oplus R^{\oplus n} =
(P \oplus R^{\oplus n}) \oplus R^{\oplus m}
$$
is free. Thus $P''$ is stably free. By symmetry we get the last of the
three cases.
\end{proof}

\end{document}
